\documentclass[10pt,a4paper]{amsart}
\usepackage[margin=19mm]{geometry}
\usepackage{amsmath,amssymb,mathtools}
\usepackage[scr=rsfs,cal=euler]{mathalfa}
\usepackage{xfrac}
\usepackage[unicode,hidelinks,bookmarks=false]{hyperref}
\newcommand{\ie}{i.e.\ }
\newcommand{\cf}{cf.\ }
\newcommand{\resp}{respectively\ }
\newcommand{\Subsection}{\subsection}
\providecommand{\nobreakdash}{\nobreak\hbox{-}}
\setlength{\emergencystretch}{2em}
\multlinegap0pt
\usepackage{tikz}
\usepackage{mathtools}
\usepackage{amssymb}
\usepackage{cancel}
\usepackage{stackengine}
\usepackage{dsfont}
\usepackage{xcolor}
\usepackage{enumerate}
\newtheorem{theorem}{Theorem}
\newcommand{\C}{\mathbb{C}}
\newcommand{\D}{\mathcal{D}}
\newcommand{\E}{\mathcal{E}}
\newcommand{\M}{\mathcal{M}}
\renewcommand{\ker}{\operatorname{Ker}}
\newcommand{\la}{\lambda}
\newcommand{\p}{\partial}
\title{Diassociative conformal algebras and their deformation cohomology}
\author{, Anupam Sahoo, Apurba Das}
\date{Source: arXiv:2609.17441v1 (CC BY 4.0)}
\begin{document}
\maketitle

\noindent\emph{Source and licence.} Diassociative conformal algebras and their deformation cohomology. Version arXiv:2609.17441v1 (CC BY 4.0), distributed by arXiv under the Creative Commons Attribution 4.0 International licence, http://creativecommons.org/licenses/by/4.0/, as recorded in the arXiv licence metadata for this exact version and shown on the abstract page https://arxiv.org/abs/2609.17441v1. Full licence notice: \texttt{licenses/arxiv-2609.17441-CC-BY-4.0.txt}.\par\smallskip

\noindent\textit{Editorial scope.} Counted unit: the article's theorem thm-ext (source0.tex lines 1472--1474). The article's complete original proof of it is delivered with it (source0.tex lines 1476--1620, one \texttt{proof} environment of 145 lines). No mathematics is written, repaired or re-derived: the statement and the proof are the article's own lines, in the article's own notation, and no retained line is substituted or re-flowed.  No further source-local context is needed: every cross-reference of every retained line resolves inside the two delivered spans.  Nothing was omitted from any retained span, so the package carries no omission record.  No retained line contains a citation, so the wrapper carries no bibliography and no bibliography entry.  The retained lines contain the article's own diagram code, which the wrapper typesets with the standard tikz packages; no image file is delivered and the package declares no asset dependency.  Editorial formatting only, in addition: an AMS \texttt{amsart} preamble that restates the article's own declarations and macros used by the retained lines, namely the environments \texttt{theorem} on separate counters, together with the standard packages those lines need.  The retained environments print in this document as Theorem 1 in order of appearance, whereas the article prints the same items with the numbering of its own full document.  The complete native source of the article as distributed in the arXiv e-print for this exact version is bundled unchanged as \texttt{sources/210598/source0.tex}.
\begin{theorem}\label{thm-ext}
Let $(\D,\dashv_{\la},\vdash_{\la})$ be a diassociative conformal algebra and $\M$ be a representation of it. Then there is a natural bijection between $\operatorname{Ext}(\D,\M)$ and $H_{cDiass}^{2}(\D,\M)$.
\end{theorem}

\begin{proof}
Let $0\longrightarrow
\M
\xrightarrow{\ i \ }
\E
\xrightarrow{\ \pi\ }
\D
\longrightarrow 0$ be a $\C[\p]$-split extension of the diassociative conformal algebra $\D$ by the representation $\M$. Since the extension is $\C[\p]$-split, we may identify $\E$ with the direct sum $\M \oplus \D$ as a $\C[\p]$-module. Under this identification, we have $i(u)=(u,0), ~ \! \pi(u,a)=a$ and $\p(u ,a)=(\p u,\p a)$, for all $u \in \M$ and $ a \in \D$. Since $\pi$ is a homomorphism of diassociative conformal algebras,
\begin{align*}
\pi ((0,a)\dashv_{\la}(0,b))=a \dashv_{\la} b ~~~ \text{ and } ~~~
\pi((0,a) \vdash_{\la} (0,b))=a \vdash_{\la} b,  \text{ for all } a, b \in \D.
\end{align*}
Therefore, we obtain a conformal sesquilinear map $f : \C[Y_2]\otimes \D^{\otimes 2 } \to \M[\la]$ such that
\begin{align*}
(0, a) \dashv_{\la}(0,b)
=
\big(f_{\la} (\tikz[baseline=-0.5ex,scale=0.20]{
  \draw (-1,1)--(0,0);
  \draw (1,1)--(0,0);
  \draw (0,0)--(0,-1);
  \draw (0.5,0.5)--(0,1);
};a, b ), \, a \dashv_{\la} b \big) ~~~ \text{ and } ~~~ (0,a) \vdash_{\la} (0,b) = \big(f_{\la} (\tikz[baseline=-0.5ex,scale=0.20]{
  \draw (-1,1)--(0,0);
  \draw (1,1)--(0,0);
  \draw (0,0)--(0,-1);
  \draw (-0.5,0.5)--(0,1);
};  a, b ), \, a \vdash_{\la} b \big).
\end{align*}
Since the representation on $\M$ induced by the extension coincides with the given representation and the $\la$-products on $\M$ are trivial, it follows that the $\la$-products of $\E$ are given by
\begin{align}
(u, a) \dashv_{\la} (v ,b) &=
\big(
a \dashv_{\la} v
+
u \dashv_{\la} b
+
f_{\la} (\tikz[baseline=-0.5ex,scale=0.20]{
  \draw (-1,1)--(0,0);
  \draw (1,1)--(0,0);
  \draw (0,0)--(0,-1);
  \draw (0.5,0.5)--(0,1);
};a, b ),
\,
a \dashv_{\la} b
\big), \label{p-1}\\
(u, a)\vdash_{\la}(v, b)
&=
(
a \vdash_{\la} v
+
u \vdash_{\la} b
+
f_{\la} (\tikz[baseline=-0.5ex,scale=0.20]{
  \draw (-1,1)--(0,0);
  \draw (1,1)--(0,0);
  \draw (0,0)--(0,-1);
  \draw (-0.5,0.5)--(0,1);
};a, b ),
\,
a \vdash_{\la} b
) \label{p-2}.
\end{align}
One may easily verify that the $\la$-products (\ref{p-1}) and (\ref{p-2}) define a diassociative conformal algebra structure on $\E$ if and only if the conformal sesquilinear map $f$ satisfies the condition $ \delta_{cDiass}(f)=0$. Equivalently, $f \in Z_{cDiass}^{2}(\D,\M)$ is a $2$-cocycle.

Next, suppose $0\longrightarrow
\M
\xrightarrow{\ i \ }
\E
\xrightarrow{\ \pi\ }
\D
\longrightarrow 0$ and $0\longrightarrow
\M
\xrightarrow{\ i' \ }
\E'
\xrightarrow{\ \pi' \ }
\D
\longrightarrow 0$
are equivalent $\C[\partial]$-split extensions, where the diassociative conformal algebra structures $(\E, \dashv_{\la}, \vdash_{\la})$ and $(\E', \dashv_{\la}',\vdash_{\la}')$ are determined by the $2$-cocycles $f$ and $f'$, respectively. Then there exists an isomorphism of diassociative conformal algebras $\varphi : \E \cong \M \oplus \D \longrightarrow \M \oplus \D \cong \E'$ such that 
\begin{align}
\pi\circ \varphi (u, a)= a ~ \text{ and } ~
\varphi (u,0)=(u,0),  \text{ for } a \in \D,  u \in \M.  \label{eq4}
\end{align}
Hence, there exists a $\C [\partial]$-linear map $g : \D \longrightarrow \M$ such that $\varphi (u, a) = (u + g (a), a)$, for all $(u, a) \in \E \cong \M \oplus \D$.
%Hence, $\xi(0,d)-(0,d) \in \ker\pi = \tau(\M)$ for every $d \in \D$ and therefore there exists a unique $\mathbb{C}$-linear map $ g: \D \rightarrow \M$ satisfying 
%\begin{align*}
%    \xi(0,d) = (g(d), d),  \text{ for all } d \in \D.
%\end{align*}
%Using the identity $\xi \p = \p \xi $, it follows that $g \in \operatorname{Hom}_{\C[\p]}(\D,\M) = C^{1}_{cDiass}(\D,\M)$. 
%Now, for any $d, d' \in \D$,
%\begin{align*}
%\xi \big((0, d) \dashv_{\la} (0, d')\big)
%=
%\xi \left( f_{\la} \left(\tikz[baseline=-0.5ex,scale=0.20]{
%  \draw (-1,1)--(0,0);
%  \draw (1,1)--(0,0);
%  \draw (0,0)--(0,-1);
 % \draw (0.5,0.5)--(0,1); };d, d' \right), \,d \dashv_{\la} d'\right) =\left(f_{\la} \left(\tikz[baseline=-0.5ex,scale=0.20]{\draw (-1,1)--(0,0); \draw (1,1)--(0,0);\draw (0,0)--(0,-1);\draw (0.5,0.5)--(0,1);};d, d' \right)+g(d \dashv_{\la} d'),\,d\dashv_{\la} d'\right),\end{align*}
%whereas
%\begin{align*}\xi(0, d)\dashv_{\la}' \xi(0,d')=(g(d), d) \dashv_{\la}' (g(d'), d')=\left(d \dashv_{\la}g(d')
%+g(d)\dashv_{\la}d'
%+f_{\la}' \left(\tikz[baseline=-0.5ex,scale=0.20]{
%  \draw (-1,1)--(0,0);
%  \draw (1,1)--(0,0);
%  \draw (0,0)--(0,-1);
%  \draw (0.5,0.5)--(0,1);
%};d, d' \right),\,
%d \dashv_{\la} d'\right).
%\end{align*}
%Since $\xi$ is a homomorphism of diassociative conformal algebras, we get 
%\begin{align*}
%  f_{\la} \left(\tikz[baseline=-0.5ex,scale=0.20]{
%  \draw (-1,1)--(0,0); \draw (1,1)--(0,0);\draw (0,0)--(0,-1);\draw (0.5,0.5)--(0,1);
%};d, d' \right) -  f_{\la}' \left(\tikz[baseline=-0.5ex,scale=0.20]{ \draw (-1,1)--(0,0); \draw (1,1)--(0,0)  \draw (0,0)--(0,-1);\draw (0.5,0.5)--(0,1);};d, d' \right) = d \dashv_{\la}g(d') - g(d \dashv_{\la} d') + g(d)\dashv_{\la}d'.\end{align*}Similarly,
%\begin{align*}
%f_{\la} \left(\tikz[baseline=-0.5ex,scale=0.20]{  \draw (-1,1)--(0,0); \draw (1,1)--(0,0);\draw (0,0)--(0,-1);\draw (-0.5,0.5)--(0,1);
%};d, d' \right) - f_{\la}' \left(\tikz[baseline=-0.5ex,scale=0.20]{
%  \draw (-1,1)--(0,0);
%  \draw (1,1)--(0,0);  \draw (0,0)--(0,-1);
%  \draw (-0.5,0.5)--(0,1);
%};d, d' \right) = d \vdash_{\la}g(d') - g(d \vdash_{\la} d') + g(d)\vdash_{\la}d'.
%\end{align*}
%Combining the above two identities, we obtain $f'-f=\delta_{cDiass}(g)$. Consequently, the cocycles $f$ and $f'$ determine the same cohomology class in $H_{cDiass}^{2}(\D, \M)$.
Since $\varphi$ is a homomorphism of diassociative conformal algebras, we have
\begin{align*}
\varphi ((0,a)\dashv_{\la}(0,b) )
=
\varphi  (0,a)\dashv_{\la}' \varphi (0,b)~
\text{ and }~
\varphi ((0,a)\vdash_{\la}(0,b) )
=
\varphi (0,a)\vdash_{\la}' \varphi (0,b ),
\end{align*}
for all $a, b \in \D$. These conditions yield $f'-f=\delta_{cDiass}(g)$, which shows that $f'$ and $f$ represent the same cohomology class in $H_{cDiass}^{2}(\D,\M)$.

\medskip

Conversely, let $f\in Z_{cDiass}^{2}(\D,\M)$ be a $2$-cocycle. Then the $\la$-products defined in \eqref{p-1} and \eqref{p-2} make $\M\oplus\D$ into a diassociative conformal algebra structure. This defines a $\mathbb{C}[\partial]$-split extension of the diassociative conformal algebra $\D$ by the representation $\M$. Next, suppose $f, f' \in  Z_{cDiass}^{2}(\D,\M)$ are cohomologous, say $ f'- f = \delta_{cDiass}(g)$ for some $g\in C_{cDiass}^{1}(\D,\M)$. Then it is easy to see that the map
\begin{align*}
\varphi: \M\oplus\D\longrightarrow\M\oplus\D  \text{ given by }
\varphi(u, a) :=(u+g(a),a), \text{ for } (u, a) \in \M \oplus \D,
\end{align*}
is an isomorphism between the $\mathbb{C}[\partial]$-split extensions corresponding to the 2-cocycles $f$ and $f'$. Hence cohomologous $2$-cocycles determine equivalent $\mathbb{C}[\partial]$-split extensions. 

Finally, the two correspondences above are inverses of each other.
\end{proof}

\end{document}
